从优化器的角度看，局部条件性的影响体现为步长选择困难。对协方差自适应后端而言，以 CMA-ES 为例，同一子问题同时包含极平坦与极陡峭方向时，能推进平坦谷底的步长会对陡峭方向过大；对陡峭方向不致发散的步长，又使平坦方向推进极慢。协方差自适应虽能逐步缓解这种失配，却需要充足采样与学习。因而在有限预算下，分组所决定的局部条件性可以实质性增加后端求解难度；现有分组研究对此讨论仍然很少。